% !TeX root = ../../MainBook.tex
\section{信号流图}\label{sec:Signal_Flow_Graph}

本节讨论多个系统进行组合时的表现，信号流图是研究此问题的有力工具。系统的\textbf{信号流图}(signal flow graph)由香农（Shannon）在1942年发明，后经梅森（J.E. Mason）推广完善，成为分析线性时不变系统的有力工具，因此又称为\textbf{Mason信号流图}或\textbf{Mason图}。

\subsection{系统的组合}\label{sub:system_comb}

信号流图可以看作是系统框图的另一种表示形式。由于系统框图不能很好地展示出系统的数学模型，本书中未涉及系统框图的概念，感兴趣的读者可以参考\cite{oppenheim1997signals}\cite{signal_systems}。

\begin{definition}
    信号流图是一种有向图，其节点表示系统的变量（输入、输出或中间变量），有向边表示变量之间的线性关系，边上的权值表示线性关系的系数。信号流图可以直观地反映系统的结构特性，并且可以通过图论的方法来分析系统的传递函数。

    信号流图相关术语汇总如下：
    \begin{itemize}
        \item \textbf{节点}：表示系统的输入/输出或中间变量，分为\textbf{输入节点（源点）}、\textbf{输出节点（阱点）}和\textbf{混合节点}；
        \item \textbf{支路}：连接两个节点的有向边，表示系统，边上的权值为系统函数，又称\textbf{增益}；
        \item \textbf{通路}：沿支路方向连接的节点序列，其中，不自交的通路称为\textbf{开通路}，从源点到阱点的开通路称为\textbf{前向通路}。将通路上各支路增益之积称为\textbf{通路增益}；
        \item \textbf{环路/闭通路}：起点和终点相同且不自交的通路。将环路上各支路增益之积称为\textbf{环路增益}；
        \item \textbf{不接触环路}：没有公共节点的两环路称为不接触环路。
    \end{itemize}
\end{definition}

信号流图与电路有一定的相似之处。两个系统“串联”时，实际上是将两个系统级联，完整系统的系统函数为两个子系统的系统函数之积，乘法交换律告诉我们\textbf{线性时不变系统总是可交换的}；两个系统“并联”时，完整系统的系统函数为两个子系统的系统函数之和（因为激励独立地经过两个系统再在输出端相加）。

\begin{figure}[H]
    \centering
    \begin{minipage}{0.45\textwidth}
        \centering
        \begin{tikzpicture}[auto, >=Stealth, node distance=1.5cm]
            \tikzset{signode/.style={circle, draw, minimum size=2mm, inner sep=0pt, line width=0.6pt}}
            \node[signode, label=left:$x(t)$] (x) {};
            \node[signode, right=of x, label=below:$w(t)$] (w) {};
            \node[signode, right=of w, label=right:$y(t)$] (y) {};
            \draw[->] (x) -- node {$H_1(s)$} (w);
            \draw[->] (w) -- node {$H_2(s)$} (y);
        \end{tikzpicture}
        \caption{串联}
    \end{minipage}
    \hfill
    \begin{minipage}{0.45\textwidth}
        \centering
        \begin{tikzpicture}[auto, >=Stealth, node distance=2.5cm]
            \tikzset{signode/.style={circle, draw, minimum size=2mm, inner sep=0pt, line width=0.6pt}}
            \node[signode, label=left:$x(t)$] (x) {};
            \node[signode, right=of x, label=right:$y(t)$] (y) {};
            \draw[->] (x) to[bend left=45] node {$H_1(s)$} (y);
            \draw[->] (x) to[bend right=45] node[swap] {$H_2(s)$} (y);
        \end{tikzpicture}
        \caption{并联}
    \end{minipage}
\end{figure}

此外，还可以引入反馈回路，根据是否有反馈回路，系统分为开环系统和闭环系统。
\begin{definition}[开环系统与闭环系统]
    如果系统取后级响应为反馈信号输入到前级，则该系统称为\textbf{闭环系统}，否则称为\textbf{开环系统}。
\end{definition}

\begin{figure}[H]
    \centering
    \begin{tikzpicture}[auto, >=Stealth, node distance=1.5cm]
        \tikzset{signode/.style={circle, draw, minimum size=2mm, inner sep=0pt, line width=0.6pt}}
        
        \node[signode, label=left:$x(t)$] (x) {};
        \node[signode, right=of x] (e) {};
        \node[signode, right=of e] (pre_y) {};
        \node[signode, right=of pre_y, label=right:$y(t)$] (y) {};
        
        \draw[->] (x) -- (e);
        \draw[->] (e) -- node {$A(s)$} (pre_y);
        \draw[->] (pre_y) -- (y);
        \draw[->] (pre_y) to[bend left=50] node {$F(s)$} (e);
    \end{tikzpicture}
    \caption{闭环系统}
\end{figure}

上图展示了一个最简单的闭环系统的信号流图。我们将$A(s)$称为\textbf{开环增益}，表示不考虑反馈回路时系统的增益，与之相对，将考虑反馈回路后的增益称为\textbf{闭环增益}。

设左侧节点的输出信号为$E(s)$，计算该系统的闭环增益$H(s)$：
\begin{align}
    Y(s)&=A(s)E(s),\nonumber\\
    E(s)&=X(s)+F(s)Y(s)\nonumber\\
    \implies Y(s)&=A(s)\sbr{X(s)+F(s)Y(s)}\nonumber\\
    \implies Y(s)&=\frac{A(s)}{1-A(s)F(s)}X(s)\nonumber\\
    \implies H(s)&=\frac{Y(s)}{X(s)}=\frac{A(s)}{1-A(s)F(s)}.
\end{align}

\begin{definition}[反馈深度与反馈类型]
    将$1-A(s)F(s)$称为\textbf{反馈深度}，它是引入反馈回路导致的增益下降的倍率。

    如果反馈深度大于$1$，则系统的输出更加稳定，但增益下降，此时反馈为\textbf{负反馈}，特别地，如果反馈深度远大于$1$，称为\textbf{深度负反馈}，此时$1-A(s)F(s)\approx -A(s)F(s)$，增益近似等于$-1/F(s)$（即反馈回路系统的逆系统），与开环增益无关。

    如果反馈深度小于$1$，则系统的输出不稳定，但增益变高，称为\textbf{正反馈}，多用于振荡器；如果反馈深度近似为$0$，系统将进入\textbf{自激振荡状态}，通常电路无法正常工作，需要消除这种状态。
\end{definition}

\subsection{信号流图的分析}\label{sub:signal_flow_graph}

\begin{example}
    设$X_1,X_2$为输入节点，$Y_1,Y_2$为输出节点，其他节点为混合节点。
    \begin{figure}[H]
        \centering
        \begin{tikzpicture}[auto, >=Stealth, node distance=1.8cm, every node/.style={font=\small}]
            % node styles: hollow smaller circles
            \tikzset{
                signode/.style={circle, draw, minimum size=2mm, inner sep=0pt, line width=0.6pt},
                lab/.style={font=\small}
            }
            % nodes (empty hollow circles)
            \node[signode] (x1) {};
            \node[signode, below=1.2cm of x1] (x2) {};
            \node[signode, right=2.2cm of x1] (m1) {};
            \node[signode, right=2.2cm of x2] (m2) {};
            \node[signode, right=2.2cm of m1] (y1) {};
            \node[signode, right=2.2cm of m2] (y2) {};
            % external labels (placed to avoid overlap)
            \node[lab, left=1mm of x1] {$X_1$};
            \node[lab, left=1mm of x2] {$X_2$};
            \node[lab, right=1mm of y1] {$Y_1$};
            \node[lab, right=1mm of y2] {$Y_2$};
            % forward path gains use H_i, feedback/return gains use G_j
            \draw[->] (x1) to node[midway, above,  inner sep=1pt] {$H_1$} (m1);
            \draw[->] (x2) to node[midway, below,  inner sep=1pt] {$H_2$} (m1);
            \draw[->] (m1) to node[midway, above,  inner sep=1pt] {$H_3$} (y1);
            \draw[->] (m2) to node[midway, below,  inner sep=1pt] {$H_4$} (y2);
            % adjusted label positions for G_5 and G_6 to avoid overlap
            \draw[->] (m1) to[bend right=8] node[midway, right,xshift=6pt, yshift=4pt,  inner sep=1pt] {$G_5$} (m2);
            \draw[->] (m2) to[bend right=22] node[midway, left, xshift=-6pt,yshift=-4pt,  inner sep=1pt] {$G_6$} (m1);
            \draw[->] (m1) edge[loop above, looseness=50] node[  inner sep=1pt] {$L_1$} (m1);
        \end{tikzpicture}
    \end{figure}
    记上、下两个混合节点对应的中间变量分别为$M_1,M_2$，则
    \begin{equation}
        \begin{aligned}
            M_1&=H_1 X_1+H_2 X_2+L_1 M_1+G_5 M_2,\\
            M_2&=G_6 M_1,\\
            Y_1&=H_3 M_1,\\
            Y_2&=H_4 M_2.
        \end{aligned}
    \end{equation}
    这是一个以$M_1,M_2,Y_1,Y_2$为未知数的线性方程组，可以解得
    \begin{equation}
        \begin{aligned}
            M_1&=\frac{H_1 X_1+H_2 X_2}{1-L_1-G_5 G_6},\\
            M_2&=\frac{G_6(H_1 X_1+H_2 X_2)}{1-L_1-G_5 G_6},\\
            Y_1&=\frac{H_3(H_1 X_1+H_2 X_2)}{1-L_1-G_5 G_6},\\
            Y_2&=\frac{H_4 G_6(H_1 X_1+H_2 X_2)}{1-L_1-G_5 G_6}.
        \end{aligned}
    \end{equation}
\end{example}

信号流图是一种有向图，$n$个节点的信号流图可以通过$n$阶的\textbf{邻接矩阵}$A$来表示，其元素$a_{ij}$表示从节点$j$到节点$i$的支路增益，如果没有支路连接这两个节点，则$a_{ij}=0$。

\begin{example}
    上例中，节点依次为$X_1,X_2,M_1,M_2,Y_1,Y_2$，则对应的邻接矩阵为
    \begin{equation}
        A=\begin{bmatrix}
            0 & 0 & 0 & 0 & 0 & 0 \\
            0 & 0 & 0 & 0 & 0 & 0 \\
            H_1 & H_2 & L_1 & G_5 & 0 & 0 \\
            0 & 0 & G_6 & 0 & 0 & 0 \\
            0 & 0 & H_3 & 0 & 0 & 0 \\
            0 & 0 & 0 & H_4 & 0 & 0
        \end{bmatrix}.
    \end{equation}
    在信号流图的语境下，源点对应的行和阱点对应的列均为$0$。有冗余的行和列，不是我们希望看到的。

    事实上，我们仅仅需要中间节点对应的邻接矩阵。由于响应$Y_1,Y_2$直接取自中间节点，只需要求解$M_1,M_2$，就可获取系统的全部信息。将信号流图表示为一个线性方程组：
    \begin{equation}
        \begin{bmatrix}
            M_1 \\
            M_2
        \end{bmatrix}=
        \begin{bmatrix}
            L_1 & G_5\\
            G_6 & 0
        \end{bmatrix}
        \begin{bmatrix}
            M_1 \\
            M_2
        \end{bmatrix}+\begin{bmatrix}
            H_1 X_1 + H_2 X_2 \\
            0
        \end{bmatrix},\quad\text{即}\ \mathbf{M}=A\mathbf{M}+\mathbf{X},
    \end{equation}
    其中
    \begin{equation}
        A=\begin{bmatrix}
            L_1 & G_5\\
            G_6 & 0
        \end{bmatrix}
    \end{equation}
    为中间节点的邻接矩阵，
    \begin{equation}
        \mathbf{M}=\begin{bmatrix}
            M_1 \\
            M_2
        \end{bmatrix},\quad\mathbf{X}=\begin{bmatrix}
            H_1 X_1 + H_2 X_2 \\
            0
        \end{bmatrix}.
    \end{equation}
    整理得
    \begin{equation}
        \begin{bmatrix}
            H_1 X_1 + H_2 X_2 \\
            0
        \end{bmatrix}=\begin{bmatrix}
            1 -L_1 & -G_5\\
            -G_6 & 1
        \end{bmatrix}\begin{bmatrix}
            M_1 \\
            M_2
        \end{bmatrix},\quad\text{即}\ \mathbf{X}=(I-A)\mathbf{M}.
    \end{equation}
    只要矩阵
    \begin{equation}
        S=I-A=\begin{bmatrix}
            1 -L_1 & -G_5\\
            -G_6 & 1
        \end{bmatrix}
    \end{equation}
    可逆，就可以求解出$M_1,M_2$，我们将这个矩阵称为\textbf{系统矩阵}(system matrix)，而将$A$称为\textbf{增益矩阵}。系统矩阵的行列式称为\textbf{特征行列式}，记为$\Delta$。在本例中，有
    \begin{equation}
        \Delta=\begin{vmatrix}
            1 -L_1 & -G_5 \\
            -G_6 & 1
        \end{vmatrix}=1-L_1-G_5 G_6.
    \end{equation}
    要使系统稳定，必须有$\Delta=1-L_1-G_5 G_6\neq 0$，否则$M_1,M_2,Y_1,Y_2\to\infty$。
\end{example}

\subsection{Mason增益公式}\label{sub:mason_gain}

本小节来考虑更一般的情况：对于有$n$个中间节点$M_1,M_2,\cdots,M_n$的系统，不妨设激励接入的节点为$M_1$（否则可用叠加原理单独分析各个激励），则系统对应的线性方程组为
\begin{equation}
    \begin{aligned}
        M_1&=a_{11} M_1 + a_{12} M_2 + \cdots + a_{1n} M_n + X,\\
        M_2&=a_{21} M_1 + a_{22} M_2 + \cdots + a_{2n} M_n,\\
        &\vdotswithin{=}\\
        M_n&=a_{n1} M_1 + a_{n2} M_2 + \cdots + a_{nn} M_n.
    \end{aligned}
\end{equation}
写成矩阵形式为
\begin{equation}
    \mathbf{M}=A\mathbf{M}+\mathbf{X},
\end{equation}
其中
\begin{equation}
    \mathbf{M}=\begin{bmatrix}
        M_1 \\
        M_2 \\
        \vdots \\
        M_n
    \end{bmatrix},\quad\mathbf{X}=\begin{bmatrix}
        X \\
        0 \\
        \vdots \\
        0
    \end{bmatrix}.
\end{equation}
在系统矩阵$S=I-A$可逆时，有
\begin{equation}
    \mathbf{M}=(I-A)^{-1}\mathbf{X}=S^{-1}\mathbf{X}.
\end{equation}
因此中间节点$M_i$可由以下公式计算：
\begin{equation}
    M_i=\frac{S^*_{1i}}{\Delta}X,
\end{equation}
其中$\Delta=\det S$为特征行列式，$S^*_{1i}$为系统矩阵$S$的去掉第$1$行、第$i$列所得的代数余子式。输出节点的响应是中间节点线性组合。以$M_n$为输出为例（否则可更换节点命名并重复此过程），可以写出\textbf{Mason增益公式}：
\begin{theorem}[Mason增益公式]
    设系统有$n$个中间节点$M_1,M_2,\cdots,M_n$，激励$X$接入节点$M_1$，响应$Y$取自节点$M_n$，则
    \begin{equation}
        H=\frac{Y}{X}=(S^{-1})_{1n}=\frac{S^*_{n1}}{\Delta},
    \end{equation}
    其中$(S^{-1})_{1n}$为系统矩阵$S$的逆矩阵的第$1$行第$n$列元素，$\Delta=\det S$为特征行列式，$S^*_{n1}$为系统矩阵$S$的去掉第$n$行、第$1$列所得的代数余子式。
\end{theorem}

根据行列式的图论展开的几何意义，Mason增益公式又可以写成如下形式：
\begin{theorem}[Mason增益公式，图论形式]\label{thm:mason_gain}
    设系统有$n$个中间节点$M_1,M_2,\cdots,M_n$，激励$X$接入节点$M_1$，响应$Y$取自节点$M_n$，则
    \begin{equation}
        H=\frac{Y}{X}=\frac{\sum_{k}g_k\Delta_k}{\Delta},
    \end{equation}
    其中$g_k$为从源点到阱点的第$k$条前向通路增益，$\Delta_k$为去掉与第$k$条前向通路接触的所有节点后的特征行列式。

    特征行列式可由以下公式计算：
    \begin{equation}
        \Delta=1-\sum_a L_a+\sum_{b,c}L_b L_c-\sum_{d,e,f}L_d L_e L_f+\cdots,
    \end{equation}
    其中$\sum_a L_a$表示所有环路增益之和，$\sum_{b,c}L_b L_c$表示所有两两不接触环路增益的乘积之和，$\sum_{d,e,f}L_d L_e L_f$表示所有三者互不接触环路增益的乘积之和，依此类推。
\end{theorem}

下面来解释为什么行列式展开具有几何意义。首先，根据行列式的公式：
\begin{equation}
    \det S=\sum_{\sigma\in S_n}\sgn(\sigma)\prod_{i=1}^{n}s_{i,\sigma(i)},
\end{equation}
其中$S_n$为所有$n$个元素的排列组成的集合，$\sgn(\sigma)$为排列$\sigma$的符号。对于$S=I-A$，对角元为$s_{ii}=1-a_{ii}$，非对角元为$s_{ij}=-a_{ij}$。

我们知道，每个置换都可以分解为不相交循环的乘积，例如，
\begin{equation}
    \sigma(1)=3,\quad \sigma(3)=4,\quad \sigma(4)=1,\quad \sigma(2)=5,\quad \sigma(5)=2
\end{equation}
可以表示为
\begin{equation}
    \sigma=(1\ 3\ 4)(2\ 5),
\end{equation}
对应中间节点组成的环路
\begin{equation}
    M_1\to M_3\to M_4\to M_1,\quad M_2\to M_5\to M_2.
\end{equation}
因此，行列式的每一项$\prod_{i=1}^{n}s_{i,\sigma(i)}$都可以看作是由若干个不接触环路增益的乘积。如果环路中出现不动点$\sigma(i)=i$，则对应的$s_{ii}=1-a_{ii}$中的$1$，表示该节点不参与任何环路；如果考虑自环，则对应$-a_{ii}$。

其次，考虑正负号的问题。对于$k$-循环，它对$\sgn(\sigma)$的贡献为$(-1)^{k-1}$，而每个增益对应一个$-a_{ij}$，因此$k$-循环对行列式的贡献为$(-1)^{k-1}(-1)^k=-1$。因此，$l$个互不接触的环路增益之积的符号应为$(-1)^l$。
\begin{example}
    考虑具有六个中间节点的信号流图，我们来计算$\sigma=(1)(2)(3\ 4\ 5\ 6)$对应的行列式项：
    \begin{equation}
        \sgn(\sigma)\prod_{i=1}^{6}s_{i,\sigma(i)}=s_{1,1}s_{2,2}s_{3,4}s_{4,5}s_{5,6}s_{6,3}=(1-a_{11})(1-a_{22})a_{34}a_{45}a_{56}a_{63}.
    \end{equation}
    将因式$(1-a_{11})(1-a_{22})$展开，得到
    \begin{equation}
        a_{34}a_{45}a_{56}a_{63}-a_{11}\cdot a_{34}a_{45}a_{56}a_{63}-a_{22}\cdot a_{34}a_{45}a_{56}a_{63}+a_{11}\cdot a_{22}\cdot a_{34}a_{45}a_{56}a_{63}.
    \end{equation}
    记环路$L:M_3\to M_4\to M_5\to M_6\to M_3$，这四项分别对应以下四种情况：
    \begin{itemize}
        \item 不考虑节点$M_1,M_2$的自环，仅考虑环路$L$，符号为$-1$；
        \item 不考虑节点$M_2$的自环，考虑节点$M_1$的自环以及环路$L$，符号为$+1$；
        \item 不考虑节点$M_1$的自环，考虑节点$M_2$的自环以及环路$L$，符号为$+1$；
        \item 考虑节点$M_1,M_2$的自环以及环路$L$，符号为$-1$。
    \end{itemize}
\end{example}

最后，对于代数余子式$S^*_{n1}=\sum_{k}g_k\Delta_k$，我们考虑将$S$按第$n$行展开：
\begin{equation}
    \det S=\Delta=\sum_{j=1}^{n}s_{nj}S^*_{nj}.
\end{equation}
这样，$S^*_{n1}$就对应$s_{n1}$的系数。行列式公式中某一项包含$s_{n1}$，说明置换$\sigma$中有$\sigma(1)=n$，即环路包含$M_1\to M_n$。将该环路的其余部分补全，就得到了一条从$M_1$到$M_n$的前向通路，其前向通路增益为$g_k$。去掉与该环路接触的节点后，剩下的部分正好对应$\Delta_k$。

注意，如果原图中不含有$M_1\to M_n$的支路，则$s_{n1}=0$，此时我们只好假想有一条“增益为$0$”的支路连接$M_1$和$M_n$，才能看出代数余子式的几何意义。因此，在Mason增益公式中，我们总是需要将前向通路考虑进来，而不管是否存在$M_1\to M_n$的支路。

\begin{example}
    考虑具有$4$个中间节点的信号流图：
    \begin{figure}[H]
        \centering
        \begin{tikzpicture}[auto, >=Stealth, node distance=2.2cm]
            % 节点（2mm空心小圆）
            \node[circle, draw, minimum size=2mm, inner sep=0pt, line width=0.6pt] (M1) {};
            \node[circle, draw, minimum size=2mm, inner sep=0pt, line width=0.6pt, right=of M1] (M2) {};
            \node[circle, draw, minimum size=2mm, inner sep=0pt, line width=0.6pt, right=of M2] (M3) {};
            \node[circle, draw, minimum size=2mm, inner sep=0pt, line width=0.6pt, right=of M3] (M4) {};
            % 节点标记，避免重叠
            \node[font=\small, below=3pt of M1] {$M_1$};
            \node[font=\small, below=3pt of M2] {$M_2$};
            \node[font=\small, below=3pt of M3] {$M_3$};
            \node[font=\small, below=3pt of M4] {$M_4$};
            % 边
            \draw[->] (M1) -- node[midway, above] {$a_{12}$} (M2);
            \draw[->] (M2) -- node[midway, above] {$a_{23}$} (M3);
            % a_{24}弯曲到M4
            \draw[->] (M2) to[out=-30,in=210,looseness=1.2] node[below right=1pt and -2pt, font=\small] {$a_{24}$} (M4);
            \draw[->] (M3) -- node[midway, above] {$a_{34}$} (M4);
            % a_{21}反馈支路弯曲
            \draw[->] (M2) to[out=150,in=30,looseness=1.3] node[above left=1pt and -2pt, font=\small] {$a_{21}$} (M1);
        \end{tikzpicture}
    \end{figure}
    仍设激励接入的节点为$M_1$，输出取自节点$M_4$，则前向通路有两条：
    \begin{equation}
        M_1\to M_2\to M_3\to M_4,\quad M_1\to M_2\to M_4.
    \end{equation}
    我们来考察前向通路与系统矩阵的对应关系。增益矩阵为
    \begin{equation}
        A=\begin{bmatrix}
            0 & a_{12} & 0 & 0 \\
            a_{21} & 0 & a_{23} & a_{24} \\
            0 & 0 & 0 & a_{34} \\
            0 & 0 & 0 & 0
        \end{bmatrix}.
    \end{equation}
    因此系统矩阵为
    \begin{equation}
        S=I-A=\begin{bmatrix}
            1 & s_{12} & 0 & 0 \\
            s_{21} & 1 & s_{23} & s_{24} \\
            0 & 0 & 1 & s_{34} \\
            0 & 0 & 0 & 1
        \end{bmatrix}.
    \end{equation}

    这个系统并没有$M_4\to M_1$的支路，但我们假想有支路$a_{41}=-s_{41}=0$，这样就产生了环路
    \begin{equation}
        M_1\to M_2\to M_3\to M_4\to M_1,\quad M_1\to M_2\to M_4\to M_1.
    \end{equation}
    前向通路增益分别为$g_1=a_{12}a_{23}a_{34}$和$g_2=a_{12}a_{24}$。

    计算特征行列式：原图中有环路$M_1\to M_2\to M_1$，环路增益为$L_1=a_{12}a_{21}$，特征行列式为$\Delta=1-L_1=1-a_{12}a_{21}$；对于两条前向通路，去掉所涉及的节点后不再有回路，因此$\Delta_1=\Delta_2=1$。根据Mason增益公式，有
    \begin{equation}
        H=\frac{Y}{X}=\frac{g_1\Delta_1+g_2\Delta_2}{\Delta}=\frac{a_{12}a_{23}a_{34}+a_{12}a_{24}}{1-a_{12}a_{21}}.
    \end{equation}
\end{example}

一般来说，我们直接用计算机分析信号流图表示的复杂系统，但在一些特定的信号流图中，使用几何意义更为方便。

\begin{example}[单入单出信号流图]
    系统的信号流图如图所示：
    \begin{figure}[H]
        \centering
        \begin{tikzpicture}[auto, >=Stealth, node distance=1.8cm, every node/.style={font=\small}]
            \tikzset{signode/.style={circle, draw, minimum size=2mm, inner sep=0pt, line width=0.6pt}, lab/.style={font=\small}}
            % nodes
            \node[signode] (X) {};
            \node[signode, right=2.2cm of X] (m1) {};
            \node[signode, below=1.2cm of m1,xshift=24mm] (m2) {};
            \node[signode, right=2.4cm of m1, yshift=12mm] (m3) {};
            \node[signode, right=2.2cm of m3] (Y) {}; 
            % labels
            \node[lab,left=1mm of X] {$X$};
            \node[lab,right=1mm of Y] {$Y$};
            \node[lab, above=0.5mm of m1] {$M_1$};
            \node[lab, below=0.5mm of m2] {$M_2$};
            \node[lab, above=0.5mm of m3] {$M_3$};
            % input connects only to m1 with unit gain
            \draw[->] (X) to node[midway, above, inner sep=1pt] {$1$} (m1);
            % output is taken from m3 with unit gain
            \draw[->] (m3) to node[midway, above, inner sep=1pt] {$1$} (Y);
            % complete directed subgraph among m1..m3 (H_ij for forward, G_ij for feedback)
            \draw[->] (m1) to[bend left=18] node[midway, above, inner sep=1pt] {$H_{12}$} (m2);
            \draw[->] (m2) to[bend left=18] node[midway, below, inner sep=1pt] {$G_{21}$} (m1);
            \draw[->] (m1) to[bend left=18] node[midway, above, inner sep=1pt,xshift=-6pt] {$H_{13}$} (m3);
            \draw[->] (m3) to[bend left=18] node[midway, below, inner sep=1pt,xshift=6pt] {$G_{31}$} (m1);
            \draw[->] (m2) to[bend left=18] node[midway, above, inner sep=1pt] {$H_{23}$} (m3);
            \draw[->] (m3) to[bend left=18] node[midway, below, inner sep=1pt] {$G_{32}$} (m2);
        \end{tikzpicture}
    \end{figure}
    \ding{172} 使用矩阵进行分析：系统矩阵为
    \begin{equation}
        \begin{bmatrix}
            1 & -G_{21} & -G_{31} \\
            -H_{12} & 1 & -G_{32} \\
            -H_{13} & -H_{23} & 1
        \end{bmatrix}.
    \end{equation}
    特征行列式为
    \begin{equation}
        \Delta=\begin{vmatrix}
            1 & -G_{21} & -G_{31} \\
            -H_{12} & 1 & -G_{32} \\
            -H_{13} & -H_{23} & 1
        \end{vmatrix}=1-G_{21}H_{12}-G_{31}H_{13}-G_{32}H_{23}-G_{21}G_{32}H_{13}-G_{31}H_{23}H_{12}.
    \end{equation}
    在$\Delta\neq 0$时，系统矩阵可逆，其逆矩阵为
    \begin{equation}
        S^{-1}=\frac{1}{\Delta}\begin{bmatrix}
            1-G_{32}H_{23} & G_{21}+G_{31}H_{23} & G_{31}+G_{21}G_{32} \\
            H_{12}+H_{13}G_{32} & 1-G_{31}H_{13} & G_{32}+G_{31}H_{12} \\
            H_{13}+H_{12}H_{23} & H_{23}+G_{21}H_{13} & 1-G_{21}H_{12}
        \end{bmatrix}.
    \end{equation}
    输出为$\displaystyle Y=M_3=\frac{1}{\Delta}(H_{13}+H_{12}H_{23})X$，因此系统增益为
    \begin{equation}
        H=\frac{Y}{X}=\frac{H_{13}+H_{12}H_{23}}{\Delta}=\frac{H_{13}+H_{12}H_{23}}{1-G_{21}H_{12}-G_{31}H_{13}-G_{32}H_{23}-G_{21}G_{32}H_{13}-G_{31}H_{23}H_{12}}.
    \end{equation}

    \ding{173} 使用Mason增益公式的几何意义：图中共有五个环路，分别为
    \begin{equation}
        \begin{aligned}
            L_1&:M_1\to M_2\to M_1,\quad \text{增益}\ L_1=G_{21}H_{12},\\
            L_2&:M_1\to M_3\to M_1,\quad \text{增益}\ L_2=G_{31}H_{13},\\
            L_3&:M_2\to M_3\to M_2,\quad \text{增益}\ L_3=G_{32}H_{23},\\
            L_4&:M_1\to M_2\to M_3\to M_1,\quad \text{增益}\ L_4=G_{21}G_{32}H_{13},\\
            L_5&:M_1\to M_3\to M_2\to M_1,\quad \text{增益}\ L_5=G_{31}H_{23}H_{12}.
        \end{aligned}
    \end{equation}
    这五个环路两两相交，因此特征行列式为
    \begin{equation}
        \Delta=1-(L_1+L_2+L_3+L_4+L_5)=1-G_{21}H_{12}-G_{31}H_{13}-G_{32}H_{23}-G_{21}G_{32}H_{13}-G_{31}H_{23}H_{12}.
    \end{equation}
    前向通路共有两条，增益分别为$g_1=H_{13}$和$g_2=H_{12}H_{23}$，去掉前向通路中的节点后没有环路，因此$\Delta_1=\Delta_2=1$。根据Mason增益公式，有
    \begin{equation}
        H=\frac{g_1\Delta_1+g_2\Delta_2}{\Delta}=\frac{H_{13}+H_{12}H_{23}}{1-G_{21}H_{12}-G_{31}H_{13}-G_{32}H_{23}-G_{21}G_{32}H_{13}-G_{31}H_{23}H_{12}}.
    \end{equation}
\end{example}

\begin{definition}[转置系统]
    将一个信号流图中的所有箭头反向，系统函数不变，所得的系统称为\textbf{转置系统}(transposed system)，其信号流图称为\textbf{对偶信号流图}(dual signal flow graph)。
\end{definition}
利用Mason增益公式，可以证明：
\begin{theorem}
    对于单入单出(SISO)系统，其转置系统的系统函数与原系统相同。
\end{theorem}
这是因为，转置系统的前向通路与原系统的前向通路一一对应，增益相同；转置系统的环路与原系统的环路一一对应，增益相同；转置系统中不接触环路的组合与原系统中不接触环路的组合一一对应，增益相同。因此，转置系统的特征行列式与原系统相同，代数余子式也相同，从而Mason增益公式给出的系统函数也相同。

\subsection{信号流图的标准型}\label{sub:standard_form}

我们同样能够根据系统函数画出信号流图，但是一个系统函数可以对应多种信号流图，例如两个系统的级联是可交换的。下面给出一种仅仅使用乘法器和积分器（系统函数为$s^{-1}$）将有理系统函数化为信号流图的标准方法，使用这种方法画出的信号流图是唯一的，并且具有相似的结构。

考虑如下信号流图：

\begin{center}
\scalebox{0.92}{%
\begin{tikzpicture}[>=Stealth, node distance=1.6cm, auto]
    % 节点：M1 为矩形，其余中间节点用直径为2的小圆圈表示，节点内部不显示编号
    \node (x) {$x(t)$};
    \node[draw, circle, minimum size=2mm, inner sep=0pt] (M1) [right=of x] {};
    \node[circle, draw, minimum size=2mm, inner sep=0pt] (M2) [right=of M1] {};
    \node[circle, draw, minimum size=2mm, inner sep=0pt] (M3) [right=of M2] {};
    \node (dots) [right=of M3] {$\cdots$};
    \node[circle, draw, minimum size=2mm, inner sep=0pt] (Mn_2) [right=of dots] {};
    \node[circle, draw, minimum size=2mm, inner sep=0pt] (Mn_1) [right=of Mn_2] {};
    \node[circle, draw, minimum size=2mm, inner sep=0pt] (Mn) [right=of Mn_1] {};
    \node (y) [right=of Mn] {$y(t)$};

    % 串联 s^{-1} 通路（链路上方标注，避免与其他箭头重叠）
    \draw[->] (x) -- (M1);
    \draw[->] (M1) -- node[above] {$s^{-1}$} (M2);
    \draw[->] (M2) -- node[above] {$s^{-1}$} (M3);
    \draw[->] (M3) -- node[above] {$s^{-1}$} (dots);
    \draw[->] (dots) -- node[below] {$s^{-1}$} (Mn_2);
    \draw[->] (Mn_2) -- node[below] {$s^{-1}$} (Mn_1);
    \draw[->] (Mn_1) -- node[below] {$b_n$} (Mn);
    \draw[->] (Mn) -- (y);

    % 直接到 Mn 的前向通路，标签放在上方并调整弯度避免重叠
    \draw[->, bend left=60] (M1) to node[above, yshift=1pt] {$b_0$} (Mn);
    \draw[->, bend left=50] (M2) to node[above, yshift=1pt] {$b_1$} (Mn);
    \draw[->, bend left=40] (M3) to node[above, yshift=1pt] {$b_2$} (Mn);
    \draw[->, bend left=30] (Mn_2) to node[above, yshift=1pt] {$b_{\,n-1}$} (Mn);

    % 反馈回到 M1（画在下方，标签置于弧线下方并微调，避免文字与箭头相交）
    \draw[->, bend left=30] (M2) to node[below, yshift=0pt] {$-a_1$} (M1);
    \draw[->, bend left=40] (M3) to node[below, yshift=-1pt] {$-a_2$} (M1);
    \draw[->, bend left=50] (Mn_2) to node[below, yshift=-1pt] {$-a_{n-1}$} (M1);
    \draw[->, bend left=60] (Mn_1) to node[below, yshift=-1pt] {$-a_{n}$} (M1);
\end{tikzpicture}%
}
\end{center}

注意到这个信号流图的所有前向通路和环路均经过节点$M_1$，去掉节点$M_1$后不再有环路，因此不含$M_1$的子图的特征行列式$\Delta_k$均为$1$。利用Mason增益公式，可以计算出系统增益为
\begin{equation}
    H(s)=\frac{Y(s)}{X(s)}=\frac{b_0 + b_1 s^{-1} + b_2 s^{-2} + \cdots + b_n s^{-n}}{1 + a_1 s^{-1} + a_2 s^{-2} + \cdots + a_n s^{-n}}.
\end{equation}

可见，对于任意有理系统函数，都可以使用上述信号流图表示。注意系数$a_k,b_k$是可以为$0$的，因此即便分子分母不是相同度数的多项式，也可以使用该信号流图表示。
